\section{Sua parte: renda em U, transformação e construção do modelo}
\textbf{Mensagem principal:} a relação com a renda é em U invertido; a resposta entra em log (Box-Cox); e o modelo foi montado em blocos, cada um testado pelo $F$ parcial.

\subsection{O que enfatizar}
\begin{itemize}
\item \textbf{Renda:} nos dados, o pico fica perto de R\$ \nRendaDecilPico; os mais violentos são os de renda média, não os mais pobres. Por isso entram renda e renda² (modelo polinomial).
\item \textbf{Box-Cox:} $\lambda$ estimado \nLambda{} (IC \nLambdaInf{} a \nLambdaSup); usamos o log, simples de interpretar. Normalidade: $W$ de \nSWy{} para \nSWln.
\item \textbf{Equação:} cada termo corresponde a uma decisão já mostrada: log da taxa (Box-Cox), renda e renda² (U invertido), os blocos A a D e o estado (variação entre regiões).
\item \textbf{Blocos:} todos entram; educação e migração dão $F$ parcial \nFblocoEd; o maior ganho de $R^2_a$ vem do estado, que no lugar da região leva o $R^2_a$ de \nRdoisAregiao{} para \nRdoisA.
\end{itemize}

\subsection{Perguntas prováveis}
\begin{description}[leftmargin=0pt]
\item[O $\lambda$ estimado não é 0, e o 0 está fora do IC; por que usar o log?] A apostila recomenda usar um valor de fácil interpretação perto do estimado; o log permite ler os efeitos em \%. Dizemos no slide que a escolha é por interpretabilidade.
\item[O que é cada termo da equação?] $r$ é o ln da renda centralizado, $\mathbf{x}$ são as outras explicativas, $D$ são as indicadoras de estado e $\varepsilon$ é o erro, suposto normal, independente e de variância constante.
\item[Por que estado e não região?] O $F$ parcial do estado além da região é \nFuf{} com \nGlUF{} gl (significativo): os estados diferem muito entre si (polícia, facções), mais do que a região capta.
\end{description}

\subsection{Transição}
``Com o modelo montado, o Willian mostra o que ele diz e como checamos as hipóteses.''
